\textbf{Phần 1. Trắc nghiệm lựa chọn}
\setcounterpageref{ex}{0}

\begin{ex} %[2D8V43]
Kết quả của phép chia $14 : 2$ là
\choice
{$6$}
{$8$}
{\True $7$}
{$5$}
\loigiai{
$14 : 2 = 7$.
\textbf{Đáp án của câu hỏi là C}
}
\end{ex}

\begin{ex} %[2D8V43]
Có $16$ cái bánh chia đều cho hai chị em. Mỗi người được bao nhiêu cái bánh?
\choice
{\True $8$ cái bánh}
{$7$ cái bánh}
{$9$ cái bánh}
{$6$ cái bánh}
\loigiai{
$16 : 2 = 8$ (cái bánh).
\textbf{Đáp án của câu hỏi là A}
}
\end{ex}

\begin{ex} %[2D8V43]
Tìm $x$ biết: $x \times 2 = 12$.
\choice
{$5$}
{$7$}
{$8$}
{\True $6$}
\loigiai{
$x = 12 : 2 = 6$.
\textbf{Đáp án của câu hỏi là D}
}
\end{ex}

\begin{ex} %[2D8V43]
Tính giá trị biểu thức: $18 : 2 + 15$.
\choice
{$22$}
{\True $24$}
{$23$}
{$25$}
\loigiai{
$18 : 2 = 9$; $9 + 15 = 24$.
\textbf{Đáp án của câu hỏi là B}
}
\end{ex}

\begin{ex} %[2D8V43]
Một đoạn dây dài $20\text{ cm}$ gập đôi lại. Chiều dài mỗi nửa đoạn dây là
\choice
{$8\text{ cm}$}
{$12\text{ cm}$}
{\True $10\text{ cm}$}
{$5\text{ cm}$}
\loigiai{
$20 : 2 = 10\text{ (cm)}$.
\textbf{Đáp án của câu hỏi là C}
}
\end{ex}

\begin{ex} %[2D8V43]
Có $18$ học sinh xếp thành $2$ hàng đều nhau. Mỗi hàng có bao nhiêu học sinh?
\choice
{\True $9$ học sinh}
{$8$ học sinh}
{$10$ học sinh}
{$7$ học sinh}
\loigiai{
$18 : 2 = 9$ (học sinh).
\textbf{Đáp án của câu hỏi là A}
}
\end{ex}

\begin{ex} %[2D8V43]
Điền số thích hợp: $16 : 2 = \dots + 3$.
\choice
{$4$}
{$6$}
{$7$}
{\True $5$}
\loigiai{
$16 : 2 = 8$. Ta có: $8 - 3 = 5$. Vậy điền số 5.
\textbf{Đáp án của câu hỏi là D}
}
\end{ex}

\begin{ex} %[2D8V43]
Bác Nam có $10$ con chim bồ câu nhốt đều vào $2$ lồng. Mỗi lồng có bao nhiêu con chim?
\choice
{$4$ con}
{\True $5$ con}
{$6$ con}
{$3$ con}
\loigiai{
$10 : 2 = 5$ (con).
\textbf{Đáp án của câu hỏi là B}
}
\end{ex}

\begin{ex} %[2D8V43]
Tìm một số biết rằng lấy số đó chia cho $2$ rồi cộng với $8$ thì được $15$.
\choice
{$12$}
{$16$}
{\True $14$}
{$18$}
\loigiai{
Tính ngược từ cuối: $15 - 8 = 7$; $7 \times 2 = 14$.
\textbf{Đáp án của câu hỏi là C}
}
\end{ex}

\begin{ex} %[2D8V43]
Thương của số lớn nhất trong bảng nhân $2$ và số $2$ là
\choice
{\True $10$}
{$9$}
{$8$}
{$20$}
\loigiai{
Số lớn nhất trong bảng nhân 2 là 20. Thương: $20 : 2 = 10$.
\textbf{Đáp án của câu hỏi là A}
}
\end{ex}

\begin{ex} %[2D8V43]
Điền dấu thích hợp: $12 : 2 + 4 \dots 18 : 2$.
\choice
{$<$}
{$=$}
{Không so sánh được}
{\True $>$}
\loigiai{
Vế trái: $6 + 4 = 10$. Vế phải: $18 : 2 = 9$. Vì $10 > 9$ nên điền dấu $>$. 
\textbf{Đáp án của câu hỏi là D}
}
\end{ex}

\begin{ex} %[2D8V43]
Một người thợ cưa một cây gỗ dài $8\text{ m}$ thành $2$ đoạn bằng nhau. Mỗi đoạn dài bao nhiêu mét?
\choice
{$3\text{ m}$}
{\True $4\text{ m}$}
{$5\text{ m}$}
{$2\text{ m}$}
\loigiai{
$8 : 2 = 4\text{ (m)}$.
\textbf{Đáp án của câu hỏi là B}
}
\end{ex}

\textbf{Phần 2. Trắc nghiệm đúng sai}
\setcounterpageref{ex}{0}

\begin{ex} %[2D8V43]
Xét tính đúng sai của các phép tính trong bảng chia 2:
\choiceTF
{\True Một số chẵn chia cho 2 luôn được thương là số tự nhiên.}
{Phép chia 15 : 2 là một phép chia hết trong bảng chia 2.}
{\True Trong bảng chia 2, số bị chia lớn nhất là 20, thương lớn nhất là 10.}
{Biểu thức 12 : 2 có kết quả bé hơn biểu thức 10 : 2.}
\loigiai{
\begin{itemize}
    \item Ý a) Đúng: Đúng vì mọi số chẵn đều chia hết cho 2.
    \item Ý b) Sai: 15 là số lẻ nên chia 2 có dư.
    \item Ý c) Đúng: Đúng.
    \item Ý d) Sai: $12 : 2 = 6 > 10 : 2 = 5$. Sai.
\end{itemize}
\begin{itemize}
    \item a) \textbf{Đúng}.
    \item b) \textbf{Sai}.
    \item c) \textbf{Đúng}.
    \item d) \textbf{Sai}.
\end{itemize}
}
\end{ex}

\begin{ex} %[2D8V43]
Mẹ có 14 quả táo chia đều cho hai anh em An và Bình:
\choiceTF
{\True Mỗi người nhận được 7 quả táo.}
{\True Số táo của An bằng số táo của Bình.}
{Nếu mẹ cho thêm mỗi người 1 quả thì mỗi người có 9 quả.}
{\True Tổng số táo của hai người là một số chẵn.}
\loigiai{
\begin{itemize}
    \item Ý a) Đúng: $14 : 2 = 7$ quả.
    \item Ý b) Đúng: Đúng vì chia đều.
    \item Ý c) Sai: Mỗi người có $7 + 1 = 8$ quả (không phải 9 quả).
    \item Ý d) Đúng: 14 là số chẵn. Đúng.
\end{itemize}
\begin{itemize}
    \item a) \textbf{Đúng}.
    \item b) \textbf{Đúng}.
    \item c) \textbf{Sai}.
    \item d) \textbf{Đúng}.
\end{itemize}
}
\end{ex}

\textbf{Phần 3. Trả lời ngắn}
\setcounterpageref{ex}{0}

\begin{ex} %[2D8V43]
Tính kết quả: 16 : 2.
\shortans{8}
\loigiai{
\textbf{Đáp số:} $8$\\
$16 : 2 = 8$.
}
\end{ex}

\begin{ex} %[2D8V43]
Có 18 viên bi chia đều cho 2 bạn. Mỗi bạn được bao nhiêu viên bi?
\shortans{9}
\loigiai{
\textbf{Đáp số:} $9$\\
$18 : 2 = 9$ viên bi.
}
\end{ex}

\begin{ex} %[2D8V43]
Tính giá trị: 12 : 2 + 18.
\shortans{24}
\loigiai{
\textbf{Đáp số:} $24$\\
$12 : 2 = 6$; $6 + 18 = 24$.
}
\end{ex}

\begin{ex} %[2D8V43]
Một đàn gà có tất cả 14 cái chân. Hỏi đàn gà có bao nhiêu con gà?
\shortans{7}
\loigiai{
\textbf{Đáp số:} $7$\\
$14 : 2 = 7$ con gà.
}
\end{ex}

\begin{ex} %[2D8V43]
Tìm x biết: x x 2 = 20.
\shortans{10}
\loigiai{
\textbf{Đáp số:} $10$\\
$x = 20 : 2 = 10$.
}
\end{ex}

\begin{ex} %[2D8V43]
Một đoạn dây dài 10 dm gập đôi lại. Hỏi mỗi nửa dài bao nhiêu đề-xi-mét?
\shortans{5}
\loigiai{
\textbf{Đáp số:} $5$\\
$10 : 2 = 5\text{ dm}$.
}
\end{ex}
